\documentclass[11pt,a4paper]{article}
\usepackage[utf8]{inputenc}
\usepackage[T1]{fontenc}
\usepackage[spanish,es-nodecimaldot]{babel}
\usepackage{lmodern,microtype,geometry,booktabs,tabularx,csquotes}
\geometry{margin=2.3cm}
\usepackage{hyperref}
\hypersetup{hidelinks}
\usepackage[backend=biber,style=authoryear,sorting=nyt,maxcitenames=2,maxbibnames=10]{biblatex}
\addbibresource{references.bib}
\setlength{\parskip}{0.22em}
\title{\textbf{Estimación transductiva del rendimiento de cebada a escala parcelaria}\\[0.4em]\large Integración de información climática, expertos locales, agrupamiento geológico y corrección por análogos\\[0.5em]\normalsize Informe técnico del Reto AgroCebada FIRA 2026}
\author{Equipo GeoCebada}
\date{19 de octubre de 2026}
\begin{document}
\maketitle
\begin{abstract}
Se documenta un modelo para 197 parcelas de cebada, con 138 rendimientos conocidos y 59 ocultos. El sistema combina un modelo climático global, un experto geográfico local, una compuerta geológica y correcciones residuales. La variante seleccionada incorpora imputación por media y corrección conservadora por tres análogos. Su RMSE en 138 exclusiones auténticas fue 0.462418 toneladas por hectárea; el procedimiento alternativo con selección interna obtuvo 0.462892. Las comparaciones son exploratorias y no miden el error de las 59 parcelas sin etiquetas.
\end{abstract}
\section{Introducción}La estimación del rendimiento agrícola a escala de parcela plantea un problema de predicción en el que confluyen información ambiental heterogénea, observaciones espaciotemporales irregulares y procesos agronómicos parcialmente conocidos. La producción de cebada responde a la disponibilidad de agua, su distribución durante el ciclo, el desarrollo fenológico, las propiedades edáficas, el relieve y las prácticas de manejo. Tales efectos no necesariamente son constantes entre localidades o campañas. La literatura ha documentado la relación entre índices espectrales y rendimiento y la sensibilidad del cultivo al estrés hídrico en etapas específicas \parencite{Mirosavljevic2018,Sharifi2020,Chanev2025,Hajkova2019,Bello2022,Mogensen1980,Zhang2015,King2024}.

El Reto AgroCebada FIRA 2026 presenta un escenario delimitado: existen 197 parcelas con covariables observables, pero solamente 138 cuentan con un rendimiento registrado y las 59 restantes requieren estimación. La unidad estadística es la parcela identificada por \texttt{ID\_POLIGONO}, no cada fila satelital. El objetivo es reconstruir \texttt{RENDIMIENTO\_T\_HA}, medido en toneladas por hectárea, para un conjunto fijo de unidades correspondientes al ciclo abril--octubre de 2025 en Hidalgo, Puebla y Tlaxcala.

Este escenario se aborda como un problema transductivo. Es admisible estudiar los atributos de las 197 parcelas mediante transformaciones que no dependen del rendimiento; en cambio, ningún ajuste supervisado puede consultar las 59 etiquetas ocultas ni la etiqueta de una parcela retenida en validación. La transducción no es sinónimo de rellenar artificialmente los rendimientos faltantes ni de entrenar sobre 197 resultados conocidos. Su fundamento metodológico aparece en el aprendizaje semisupervisado y la regularización geométrica \parencite{Chapelle2006,Belkin2006,Zhou2004}.

La evidencia disponible comprende observaciones BASIC de Sentinel y Landsat, observaciones PRO de Planet, climatología y precipitación CHIRPS, productos topográficos, información de suelo y recursos hídricos. Las señales tienen distintas resoluciones, escalas, fechas y mecanismos de ausencia. Por ello, una cobertura extensa de imágenes no equivale a contar con una muestra supervisada grande; los 138 rendimientos siguen siendo la restricción estadística central \parencite{Funk2015CHIRPS,Poggio2021SoilGrids,INEGICEM,FAOWaPOR,SIAPDatos}.

La modelación requiere equilibrar flexibilidad e interpretabilidad. Los modelos locales reconocen heterogeneidad territorial, mientras las regresiones regularizadas disminuyen sensibilidad a covariables correlacionadas y tamaños muestrales limitados. Las comparaciones se realizaron mediante validación dejando fuera una parcela, con especial cuidado de que la selección de parámetros ocurra sin usar su rendimiento. No obstante, explorar muchas arquitecturas sobre las mismas 138 etiquetas puede introducir optimismo de selección \parencite{Stone1974,HoerlKennard1970}.

Se adoptó una arquitectura denominada Reference26, compuesta por un núcleo climático global, un experto geográfico local, una compuerta de composición geológica y una corrección residual. Posteriormente se evaluaron estrategias de imputación, mezclas con otros modelos y una corrección final mediante análogos de 26 atributos. Este documento explica tanto los mecanismos retenidos como las alternativas consideradas, detalla los resultados de las exclusiones auténticas y distingue las estimaciones de error observadas de la precisión todavía desconocida en las 59 parcelas objetivo.

\section{Justificación}La falta de rendimientos medidos de forma homogénea limita la caracterización productiva a escala parcelaria. La precipitación, los indicadores de vigor vegetal y las capas topográficas aportan evidencia complementaria, aunque no constituyen observaciones directas de cosecha ni identifican por sí solos las causas del rendimiento \parencite{Sharifi2020,Holland2021,Bokusheva2016,Serban2025}.

Con únicamente 138 etiquetas reales, aumentar sin restricciones la complejidad de un predictor puede producir ajustes inestables y aparentes mejoras que no se repiten en parcelas excluidas. Se justifica, por tanto, una arquitectura con componentes regularizados, vecindarios controlados y comparaciones emparejadas. La metodología prioriza el error cuadrático medio de predicción medido mediante LOOCV, pero conserva MAE, sesgo, errores extremos y diagnósticos espaciales para no ocultar defectos del modelo.

La reconstrucción de 59 rendimientos se plantea como una herramienta de apoyo técnico, no como sustitución automática de mediciones de campo o de verificación independiente. La transparencia sobre imputación, selección de modelos e incertidumbre es necesaria para evitar decisiones basadas en una precisión exagerada.

\section{Objetivos}\textbf{Objetivo general.} Construir y documentar un modelo reproducible para estimar el rendimiento de las 59 parcelas objetivo del Reto AgroCebada FIRA 2026 a partir de atributos agroambientales conocidos y de 138 rendimientos observados, con énfasis en precisión predictiva, trazabilidad y prevención de filtración de información.

\textbf{Objetivos específicos.} Primero, integrar datos satelitales, climáticos, topográficos y geológicos en registros por parcela, respetando ventanas temporales, unidades y procedencia de cada sensor. Segundo, contrastar alternativas de tratamiento de valores faltantes, como mediana, media, imputación espacial, temporal y espaciotemporal, sin utilizar el rendimiento de la parcela retenida. Tercero, comparar regresiones globales regularizadas, expertos geográficos y compuertas basadas en perfiles de composición geológica, incluyendo el examen de K-means y DBSCAN. Cuarto, investigar la influencia del número de vecinos y la penalización Ridge sobre conjuntos de validación idénticos. Quinto, estudiar correcciones residuales y ajustes conservadores por análogos con controles explícitos de soporte. Sexto, reportar RMSE, MAE, sesgo, coeficiente de determinación predictivo, sensibilidad a parcelas extremas y autocorrelación espacial. Finalmente, separar las inferencias válidas sobre las 138 parcelas conocidas de los resultados todavía no verificables para las 59 parcelas sin etiqueta.

\section{Metodología}\subsection*{Construcción de la evidencia y validación}
Se emplearon 197 parcelas y un conjunto supervisado de 138 rendimientos. BASIC reúne observaciones multitemporales Sentinel y Landsat; PRO incorpora observaciones Planet. Se agregaron indicadores del ciclo de 2025, antecedentes de vegetación, lluvia, topografía y geología con comprobaciones de identificadores y trazabilidad. CHIRPS, SoilGrids y WaPOR se tratan como fuentes indirectas con escala propia \parencite{Funk2015CHIRPS,Poggio2021SoilGrids,FAOWaPOR,Dhonthi2024}. La validación consistió en 138 exclusiones sucesivas de una parcela, utilizando las otras 137 etiquetas para cada ajuste. Las 59 etiquetas ocultas permanecieron inaccesibles. Una transformación sin dependencia del rendimiento puede consultar las 197 covariables; toda operación ajustada con rendimientos utiliza únicamente etiquetas disponibles en su entrenamiento \parencite{Stone1974,Chapelle2006}.

\subsection*{Estrategias de imputación}
Se consideraron mediana y media aritmética calculadas dentro de cada conjunto de entrenamiento, así como reconstrucción espacial, temporal y espaciotemporal de predictores faltantes. Las estrategias temporales respetaron fechas, sensores y calidad de observación; las espaciales exploraron parcelas próximas y contexto territorial. No se confundió completar una señal espectral con estimar el rendimiento. La versión final emplea media porque obtuvo un RMSE ligeramente inferior al de la mediana en una comparación de 138 exclusiones. No se atribuye superioridad frente a interpolaciones alternativas sin su tabla final auditada ni se sostiene que una diferencia mínima sea significativa.

\subsection*{Predicción climática y experto local}
El componente C--T3 aplica Elastic Net sobre cuatro entradas: lluvia total del ciclo, fracción correspondiente a meses tempranos, concentración mensual de precipitación y NDVI Sentinel--2 en mayo de 2025 \parencite{ZouHastie2005}. El experto local usa siete atributos distintos que resumen nivel y variabilidad histórica de NDVI, variación fenológica observada por Planet, estadísticos de NDWI y pendiente. Para cada parcela selecciona los 24 vecinos geográficamente más cercanos con rendimiento conocido, ajusta una regresión Ridge ponderada inversamente por distancia y conserva penalización 30 \parencite{HoerlKennard1970}. Un ajuste residual interno Elastic Net, basado en predicciones honestas del experto, complementa la estimación local.

\subsection*{Ponderación geológica y residuales}
Las diez proporciones de composición geológica, transformadas mediante raíz cuadrada y centrado, alimentan un K-means de cuatro u ocho grupos; la alternativa se elige dentro del entrenamiento. Cada grupo recibe un peso regularizado que mezcla las predicciones global y local. Las coordenadas sirven para elegir vecinos del experto y no se confunden con las variables de agrupamiento geológico \parencite{MacQueen1967}. Se ensayó DBSCAN como alternativa de agrupamiento por densidad, con manejo explícito de ruido y asignación de parcelas nuevas, sin cambiar los demás componentes \parencite{Ester1996DBSCAN}. El corrector R3 utiliza tres indicadores de estructura histórica del NDVI, persistencia de heterogeneidad y pendiente mínima; es regularizado, se desactiva si faltan indicadores y está acotado. En conjunto, los 26 atributos se distribuyen en cuatro globales, siete locales, dos coordenadas, diez geológicos y tres residuales.

\subsection*{Corrección final y elección interna}
Se examinó la cercanía entre parcelas en 26 atributos escalados robustamente y, alternativamente, en ocho componentes principales. La regla seleccionada considera los tres análogos etiquetados más próximos, ponderados inversamente por su distancia en atributos, y corrige parcialmente la predicción base con un factor de mezcla de 0.25. Una compuerta reduce el cambio cuando la cercanía o la concordancia de rendimientos vecinos indica soporte débil; el ajuste absoluto se limita a 0.5 toneladas por hectárea. La elección entre cinco recetas finalistas y la ausencia de corrección se analizó también mediante cuatro particiones internas en cada exclusión externa. Aunque este anidamiento evita que la etiqueta del caso retenido intervenga directamente en la elección interna, los candidatos habían sido preseleccionados tras explorar esas mismas 138 etiquetas.

\section{Resultados y discusión}\subsection*{Imputación, regularización y vecinos}
La imputación mediante media aritmética produjo un RMSE de 0.468138 toneladas por hectárea, frente a 0.4681855 para la mediana del Reference26 original. La ganancia absoluta fue apenas 0.000048 toneladas por hectárea, aproximadamente, y favoreció exclusivamente el criterio de error cuadrático. El MAE aumentó de 0.333673 con mediana a 0.333786 con media. La elección de la media fue, por tanto, una decisión operacional basada en el mínimo RMSE observado, sin evidencia de una superioridad estadística robusta. Se implementaron procedimientos de imputación espacial, temporal y espaciotemporal, pero no se dispone en la documentación pública consultada de sus cuatro RMSE completos y auditados: sería incorrecto declarar que todas las variantes interpoladas fueron peores.

Una búsqueda reducida contrastó 25 pares de números de vecinos y valores de penalización Ridge sobre las primeras 12 parcelas retenidas; después se compararon ocho configuraciones sobre las primeras 46 y tres sobre las 138 parcelas. En la última etapa, 24 vecinos y penalización 30 obtuvieron RMSE de 0.468138, frente a 0.470112 para 20 vecinos y penalización 30, y 0.474493 para 12 vecinos y penalización 3.2. El orden relativo cambió conforme creció el número de exclusiones. Este resultado muestra el peligro de convertir una ventaja preliminar sobre pocas parcelas en una selección definitiva.

\subsection*{Agrupamiento geológico: límites de la comparación}
La compuerta geológica retenida utiliza K-means para aprender pesos de mezcla entre el predictor climático global y el experto local. La alternativa DBSCAN fue examinada sustituyendo únicamente la regla de agrupamiento, conservando las diez variables geológicas y el resto de la arquitectura \parencite{MacQueen1967,Ester1996DBSCAN}. En las primeras 46 parcelas del estudio de compuertas, DBSCAN obtuvo RMSE de 0.545343, mientras el K-means original registró 0.547536. En ese subconjunto \emph{DBSCAN fue mejor en RMSE}. Las 46 parcelas son una selección fija y no equivalen a una confirmación completa sobre 138 exclusiones. Por ello se retuvo la arquitectura Reference26 con K-means previamente validada y se dejó DBSCAN como comparación exploratoria pendiente de confirmación completa. Esta justificación se refiere a la continuidad y trazabilidad de la arquitectura, no a que K-means haya ganado esa prueba.

\subsection*{Comparación de arquitecturas y selección de la corrección}
Reference26 con media obtuvo RMSE de 0.468138, MAE de 0.333786 y coeficiente de determinación predictivo de 0.695292, sobre 138 exclusiones genuinas. Compact14, arquitectura de 14 entradas y sin compuerta geológica, produjo RMSE de 0.469673 y MAE de 0.333044. Una mezcla global aproximada de 70.5\% Reference26 con 29.5\% Compact14 alcanzó RMSE de 0.467811. Su mejora frente al primero fue de solo 0.000327 toneladas por hectárea y la ventaja desapareció al excluir la parcela con mayor contribución favorable. El ensamble no se promovió como una ganancia robusta.

\begin{table}[htbp]
\centering
\small
\caption{Resultados de 138 exclusiones reales: errores en t/ha; el R\textsuperscript{2} es predictivo.}
\label{tab:resultados}
\begin{tabular}{@{}lrrr@{}}
\toprule
Modelo & RMSE & MAE & R\textsuperscript{2}\\
\midrule
Reference26 original, mediana & 0.468186 & 0.333673 & aprox.\ 0.6952\\
Reference26, media & 0.468138 & 0.333786 & 0.695292\\
Compact14 & 0.469673 & 0.333044 & 0.693291\\
Mezcla Reference26--Compact14 & 0.467811 & 0.332489 & 0.695717\\
Análogos, receta fija elegida & \textbf{0.462418} & 0.332019 & \textbf{0.702693}\\
Selección anidada entre finalistas & 0.462892 & \textbf{0.331956} & 0.702083\\
\bottomrule
\end{tabular}
\end{table}

La corrección fija por análogos robustos de 26 entradas y tres vecinos, con mezcla 0.25 y compuerta de soporte, obtuvo el menor RMSE observado: 0.462418, una reducción de 0.005720 toneladas por hectárea, aproximadamente 1.22\% respecto a Reference26 con media. Alteró las 138 predicciones, mejoró el error cuadrático de 65 y lo empeoró en 73; el mayor cambio absoluto de predicción fue de 0.108689 toneladas por hectárea. Este resultado muestra una mejora agregada pequeña, sin mejora individual uniforme. La política de selección mediante cuatro particiones internas eligió esa receta en 128 de las 138 exclusiones, una alternativa PCA de ocho componentes en cinco y ninguna corrección en cinco. Su RMSE fue 0.462892, con MAE de 0.331956 y R\textsuperscript{2} de 0.702083.

\subsection*{Residuos, dependencia espacial e incertidumbre}
El modelo corregido fijo presentó sesgo medio de -0.004103 toneladas por hectárea, error absoluto mediano de 0.244838, percentil 90 del error absoluto de 0.638737 y mayor error absoluto de 2.115301. La correlación de Pearson entre observado y predicho fue 0.838282. Los dos casos más difíciles continuaron siendo influyentes; por ello el valor agregado de RMSE no describe de forma exhaustiva el riesgo de errores extremos.

El índice de Moran de los residuos firmados fue -0.065382 para cuatro vecinos geográficos y -0.079994 para ocho, con valores bilaterales nominales 0.2715 y 0.0580, respectivamente. Para doce vecinos se observó un valor nominal de 0.0470, que no sobrevivió la corrección de Holm sobre la familia completa de contrastes espaciales: el valor ajustado fue 1.0. En consecuencia, no hay fundamento para afirmar dependencia residual significativa después del ajuste múltiple, pero tampoco para demostrar independencia geográfica \parencite{Moran1950,Geary1954,Holm1979,BenjaminiHochberg1995}.

Las pruebas de normalidad, heterocedasticidad y calibración, junto con el remuestreo por grupos territoriales, se interpretaron como diagnósticos condicionados a pronósticos previamente obtenidos \parencite{BreuschPagan1979,White1980,EfronTibshirani1993,Tobler1970}. No se asignó un R\textsuperscript{2} ajustado convencional ni valores de significación a inexistentes coeficientes globales del predictor compuesto. Finalmente, las variantes se eligieron después de muchas exploraciones usando las mismas 138 etiquetas; el anidamiento posterior protege cada exclusión frente a su selección interna, pero no deshace el sesgo de búsqueda histórica. El RMSE de las 59 parcelas con rendimiento oculto sigue siendo desconocido.

\section{Uso de la ciencia de datos en la inclusión financiera del sector agropecuario}La ciencia de datos puede ampliar el acceso a información productiva allí donde los registros de cosecha son incompletos o costosos de obtener. Las estimaciones parcelarias de rendimiento permiten complementar evaluaciones de capacidad productiva, planificación de flujos de efectivo, seguimiento territorial del riesgo y selección de necesidades de asistencia técnica. La observación satelital aporta indicadores comparables de vegetación, sequía y productividad hídrica cuya obtención no exige visitar simultáneamente todas las parcelas \parencite{Bokusheva2016,Dhonthi2024,FAOWaPOR,Serban2025}.

En el análisis financiero, sin embargo, rendimiento no equivale a rentabilidad, ingreso neto ni solvencia. Para valorar la capacidad de pago deben considerarse superficie realmente cosechada, precio de venta, costos de producción, exposición climática, otros ingresos y compromisos crediticios. Tampoco es lícito extrapolar sin reservas los indicadores de 138 parcelas a todos los productores de Hidalgo, Puebla y Tlaxcala; sus poblaciones, sistemas de producción y mecanismos de selección pueden ser diferentes de los contenidos en las estadísticas oficiales \parencite{FIRA2026AgroCebada,SIAPDatos}.

Una aplicación responsable requeriría conservar procedencia y fecha de cada predictor, revisar calidad de las observaciones, comprobar estabilidad territorial, cuantificar incertidumbre con información independiente y habilitar corrección humana ante situaciones atípicas. Las coordenadas y los rendimientos individuales deben protegerse por su sensibilidad comercial. Las estimaciones no deberían utilizarse como regla automática y exclusiva para aprobar o denegar un crédito, pues podrían desfavorecer sistemáticamente a parcelas o grupos insuficientemente representados.

En consecuencia, la principal contribución a la inclusión financiera no es sustituir la evaluación agronómica por una cifra algorítmica, sino ofrecer un insumo adicional, trazable y sometido a auditoría. La reducción de asimetrías informativas solamente genera valor sostenible cuando se acompaña de mecanismos que reconozcan explícitamente el error de predicción, el contexto productivo y las limitaciones de transferencia a otros ciclos.

\section{Conclusiones}La investigación desarrolló una estrategia de estimación del rendimiento de cebada para 59 parcelas sin etiqueta mediante información de otras 138 parcelas observadas y un conjunto de covariables ambientales disponible para las 197 unidades. La arquitectura elegida no corresponde a una regresión aislada: reúne un componente climático Elastic Net, un experto Ridge geográfico, una compuerta de composición geológica y una corrección residual restringida. 

La imputación por media se conservó tras registrar un RMSE ligeramente menor al de la mediana: 0.468138 frente a 0.468186 toneladas por hectárea. La diferencia fue mínima y no se tradujo en una reducción del MAE.  La búsqueda escalonada de parámetros también mostró que la clasificación de candidatos cambia al aumentar el número de parcelas excluidas; al completar las 138 exclusiones, la combinación de 24 vecinos y penalización Ridge 30 resultó la mejor de las finalistas examinadas.

En materia de agrupamiento, K-means integra la arquitectura Reference26 sometida a evaluación completa. No obstante, la prueba documentada sobre 46 parcelas favoreció a DBSCAN en RMSE. Por ello no cabe defender la compuerta elegida alegando una superioridad empírica universal frente a DBSCAN. Mantener la arquitectura con K-means responde a la procedencia y verificabilidad del modelo final, mientras la variante de densidad exige una comparación global homogénea antes de modificar el sistema.

La última corrección utiliza tres parcelas análogas en el espacio robusto de 26 atributos, con un factor parcial de mezcla y una compuerta de soporte para evitar cambios excesivos. La regla fija seleccionada obtuvo un RMSE de 0.462418 toneladas por hectárea y un coeficiente de determinación predictivo de 0.702693. Una política de selección interna obtuvo RMSE de 0.462892. La magnitud de mejora fue modesta, no benefició a todas las parcelas y no eliminó los errores extremos.

 Los 138 rendimientos se consultaron repetidamente durante la exploración de familias, parámetros y correcciones. En consecuencia, una parte de la aparente ventaja puede proceder de selección adaptativa; incluso la política posteriormente anidada está condicionada a candidatos ya cribados. Tampoco se conoce el rendimiento real de las 59 parcelas objetivo, por lo que sería incorrecto atribuirles las métricas de las exclusiones de entrenamiento.

Como conclusión, el sistema aporta una reconstrucción transductiva técnicamente documentada, con decisiones verificables y criterios explícitos para excluir etiquetas de evaluación. Para una aplicación financiera o productiva fuera del reto harían falta campañas independientes, calibración de incertidumbre y análisis de estabilidad por territorio. 

\clearpage
\nocite{*}
\printbibliography[heading=bibintoc,title={Bibliografía}]
\clearpage
\appendix\section{Anexos}\small
\subsection*{A. Matriz resumida de decisiones}
\begin{tabularx}{\linewidth}{@{}p{0.25\linewidth}X@{}}
\toprule
Decisión & Evidencia / límite\\
\midrule
Media frente a mediana & RMSE 0.468138 frente a 0.468186; el MAE con media aumentó ligeramente.\\
Imputación espacial o temporal & Alternativas estudiadas; sin clasificación global de RMSE verificada para 138 exclusiones, no se afirma derrota numérica.\\
Ridge: 24 vecinos, penalización 30 & Mejor de tres finalistas con media sobre las 138 exclusiones; cribado previo adaptativo.\\
K-means geológico & Parte de Reference26 completo; DBSCAN obtuvo menor RMSE en las 46 parcelas comparadas.\\
Compact14 y mezcla global & RMSE 0.469673 y 0.467811, respectivamente; mejora de la mezcla frágil ante la exclusión de un caso.\\
Corrección por tres análogos & RMSE fijo 0.462418; política de elección interna 0.462892; ambas sujetas a búsqueda previa.\\
\bottomrule
\end{tabularx}
\subsection*{B. Procedencia y reproducibilidad}
Repositorio: \url{https://github.com/heri-espino/GeoC}. Arquitectura original: \url{https://github.com/heri-espino/todoseboso}. Documentos de auditoría: \texttt{docs/REFERENCE26\_IMPUTATION\_LOOCV.md}; \texttt{docs/PEER\_LOOCV\_RECORDED\_RESULTS.md}; \texttt{docs/REFERENCE26\_TOP5\_NESTED\_LOOCV.md}; \texttt{docs/REFERENCE26\_FINAL59\_SELECTED.md}; \texttt{docs/REFERENCE26\_SELECTED\_TAIL\_STATISTICAL\_ANALYSIS.md}. Se documentan 138 exclusiones con 137 rendimientos de entrenamiento por ajuste. Los archivos privados por parcela y los 59 rendimientos de referencia no se publican. La selección definitiva de 59 objetivos debe verificarse contra el manifiesto generado en el entorno de cómputo y las bases oficiales.
\normalsize

\end{document}
